\section{Proof of Theorem~\ref{thm:per-client}}
\label{app:thm1}

We restate the theorem and give the full proof, organised as four
lemmas plus an assembly step. The proof uses standard machinery:
chain rule of mutual information, Mironov's RDP for the Gaussian
mechanism~\cite{mironov2017rdp}, and the Cuff-Yu max-KL to MI-DP
conversion~\cite{cuff2016mi}.

\setcounter{theorem}{0}
\renewcommand{\thetheorem}{\thesection.\arabic{theorem}}

\begin{theorem}[restated]
Under (IA1)--(IA2), for any deterministic adversary strategy $f$,
any prior $\mathbb{P} \in \FAM_{\TOPO,\omega}$, and any client $i$:
\[
I_{\mathbb{P}}\big(p_i;\,\hat p_i \;\big|\; \TOPO,\omega,\{\sigma_j\}\big)
\;\leq\;
\frac{\Tmax C^2}{2\sigma_i^2 |B|^2}
\;+\;
\ell_i^\circ.
\]
\end{theorem}

The proof decomposes into four lemmas. Lemmas~\ref{lem:chain}
and~\ref{lem:lateral} isolate the lateral-leakage term via a chain-rule
inequality; Lemmas~\ref{lem:perround} and~\ref{lem:composition} bound
the mechanism term using DP-SGD's Gaussian noise structure.

\begin{lemma}[Chain-rule decomposition]
\label{lem:chain}
For any random variables $X, Y, Z$:
\[
I(X;Y) \;\leq\; I(X;Z) + I(X;Y \mid Z).
\]
\end{lemma}
\begin{proof}
The conditioning identity gives
\(I(X;Y) = I(X;Z) + I(X;Y \mid Z) - I(X;Z \mid Y)\); the third term is
non-negative.
\end{proof}

\begin{lemma}[Per-round Gaussian-mechanism KL bound]
\label{lem:perround}
Under (IA1)--(IA2), for each round $t$ and any $D_{-i}$,
\[
\KL\Big(
\mathbb{P}_{\theta_i^{(t)} \mid D_i, \theta^{(t-1)}, D_{-i}}
\;\Big\|\;
\mathbb{P}_{\theta_i^{(t)} \mid D_i', \theta^{(t-1)}, D_{-i}}
\Big)
\;\leq\; \frac{C^2}{2\sigma_i^2 |B|^2},
\]
for any adjacent databases $D_i, D_i'$. Consequently, by
Cuff-Yu~\cite{cuff2016mi},
\[
I\big(D_i;\, \theta_i^{(t)} \;\big|\; \theta^{(t-1)}, D_{-i}\big)
\;\leq\; \frac{C^2}{2\sigma_i^2 |B|^2}.
\]
\end{lemma}
\begin{proof}
Conditional on $\theta^{(t-1)}$ and $D_{-i}$, the round-$t$ local update
\eqref{eq:dpsgd} is the Gaussian mechanism with sensitivity $C/|B|$
applied to clipped per-record gradients of $D_i$. By Mironov's
analysis~\cite{mironov2017rdp} and the standard DP-SGD machinery~\cite{abadi2016dpsgd},
the KL divergence between adjacent mechanism outputs is bounded by
$\Delta^2/(2s^2)$ with $\Delta = C/|B|$ and $s = \sigma_i C/|B|$, giving
$C^2/(2\sigma_i^2 |B|^2)$. Cuff and Yu's max-KL bound on
adjacent databases implies MI-DP with the same parameter~\cite{cuff2016mi}.
\end{proof}

\begin{lemma}[Sequential composition]
\label{lem:composition}
Under (IA1)--(IA2), sequential composition of the per-round bound
across all $\Tmax$ rounds gives:
\[
I\big(D_i;\, \Theta_i \;\big|\; D_{-i}\big)
\;\leq\;
\sum_{t=1}^{\Tmax} I\big(D_i;\, \theta_i^{(t)} \;\big|\; \theta^{(t-1)}, D_{-i}\big)
\;\leq\;
\frac{\Tmax C^2}{2\sigma_i^2 |B|^2}.
\]
\end{lemma}
\begin{proof}
The MI chain rule gives
\(I(D_i;\Theta_i \mid D_{-i}) = \sum_t I(D_i; \theta_i^{(t)} \mid \theta_i^{(<t)}, D_{-i})\).
Conditioning on $\theta^{(t-1)}$ subsumes conditioning on
$\theta_i^{(<t)}$ together with $D_{-i}$: the public global state
$\theta^{(t-1)}$ is determined by the past observed $\theta_i^{(<t)}$ and
$\Theta_{-i}^{(<t)}$, and the latter depends only on $D_{-i}$ and noise
which is independent of $D_i$ by (IA1). The
inequality
\(I(D_i; \theta_i^{(t)} \mid \theta_i^{(<t)}, D_{-i})
   \leq I(D_i; \theta_i^{(t)} \mid \theta^{(t-1)}, D_{-i})\)
follows because the larger conditioning set $\theta^{(t-1)}$ provides
strictly more information that is independent of $D_i$, so removes more
shared randomness from the MI between $D_i$ and the round-$t$ update.
Apply Lemma~\ref{lem:perround} per round and sum.
\end{proof}

\begin{lemma}[Lateral-leakage bound]
\label{lem:lateral}
Under (IA2),
\(I(p_i;\, D_{-i}) \leq \ell_i^\circ\).
\end{lemma}
\begin{proof}
Direct from the definition of $\ell_i^\circ$ as the supremum over
priors in $\FAM_{\TOPO,\omega}$ (Definition~\ref{def:leverage}).
\end{proof}

\paragraph{Assembling the theorem.}
Combining the lemmas:
\begin{align*}
I(p_i;\hat p_i \mid \cdot)
&\leq I(p_i;\Theta \mid \cdot)
  && \text{(data processing on }f\text{)} \\
&\leq I(p_i;D_{-i}) + I(p_i;\Theta \mid D_{-i})
  && \text{(Lemma~\ref{lem:chain})} \\
&\leq \ell_i^\circ + I(p_i;\Theta_i \mid D_{-i})
  && \text{(Lemma~\ref{lem:lateral} + (IA1) for }\Theta_{-i}\text{)} \\
&\leq \ell_i^\circ + I(D_i;\Theta_i \mid D_{-i})
  && \text{(data processing on } D_i \to p_i \text{)} \\
&\leq \ell_i^\circ + \frac{\Tmax C^2}{2\sigma_i^2 |B|^2}
  && \text{(Lemma~\ref{lem:composition}).}
\end{align*}
The conditioning on $\TOPO, \omega, \{\sigma_j\}$ is preserved
throughout because these quantities are constants of the deployment
(not random variables in our model). \qed

\section{Proof of Theorem~\ref{thm:allocation}}
\label{app:thm2}

The optimisation problem is
\(\min \max_i [a/\sigma_i^2 + \ell_i^\circ]\) subject to
$\sum_i \sigma_i^2 \leq U, \sigma_i^2 > 0$. The objective is
non-smooth, so we reformulate via a slack variable $K$ to obtain a
convex programme amenable to KKT analysis:
\begin{equation}
\min_{K, \{\sigma_i^2 > 0\}} K
\quad\text{s.t.}\quad
a/\sigma_i^2 + \ell_i^\circ \leq K \;\forall i,
\quad \sum_i \sigma_i^2 \leq U.
\label{eq:opt-reform}
\end{equation}
The objective is linear in $K$ and the constraints are convex (since
$1/\sigma^2$ is convex on $\sigma^2>0$ and the budget constraint is
linear). Slater's condition holds at the strictly feasible interior
point $\sigma_i^2 = U/(2n)$, and the per-client constraint is satisfied
at $K = 2na/U + \max_i \ell_i^\circ$. Strong duality therefore holds
and the KKT conditions are necessary and sufficient.

\begin{lemma}[KKT optimality]
\label{lem:kkt}
The solution to \eqref{eq:opt-reform} satisfies
\(\sigma_i^{*\,2} = a / (K^\star - \ell_i^\circ)\) for all $i$, with
\(K^\star\) the unique value satisfying
\(\sum_i a/(K^\star - \ell_i^\circ) = U\) and \(K^\star > \max_i \ell_i^\circ\).
\end{lemma}
\begin{proof}
Lagrangian:
\(\mathcal{L} = K + \sum_i \mu_i(a/\sigma_i^2 + \ell_i^\circ - K)
                + \lambda(\sum_i \sigma_i^2 - U)\).
Stationarity in $K$: $1 - \sum_i \mu_i = 0$ so $\sum_i \mu_i = 1$.
Stationarity in $\sigma_i^2$:
$-\mu_i a /\sigma_i^4 + \lambda = 0$ giving
$\sigma_i^{*\,2} = \sqrt{\mu_i a / \lambda}$. Since $\sigma_i^2 = 0$ would
require $\mu_i = 0$ which forces $a/\sigma_i^2 + \ell_i^\circ = +\infty
\not\leq K$, every per-client constraint must be active at the optimum:
$a/\sigma_i^{*\,2} + \ell_i^\circ = K^\star$ for all $i$, giving the
closed form. Uniqueness of $K^\star$: define
\(g(K) = \sum_i a/(K - \ell_i^\circ)\) on $K > \max_i \ell_i^\circ$.
This is continuous and strictly decreasing, with
\(g(K) \to \infty\) as $K \downarrow \max_i \ell_i^\circ$ and
\(g(K) \to 0\) as $K \to \infty$. By the intermediate value theorem
there is a unique $K^\star$ with $g(K^\star) = U$ for any $U > 0$.
\end{proof}

\begin{corollary}[Strict improvement, restated]
\label{cor:strict-restated}
Let $K_{\mathrm{uniform}} = an/U + \max_i \ell_i^\circ$, the worst-case
bound at uniform allocation $\sigma_i^2 = U/n$. Then
\(K^\star \leq K_{\mathrm{uniform}}\) with equality iff all
$\ell_i^\circ$ are equal.
\end{corollary}
\begin{proof}
Uniform allocation is feasible for \eqref{eq:opt-reform} with
objective $K_{\mathrm{uniform}}$, so $K^\star \leq K_{\mathrm{uniform}}$.
Equality requires uniform allocation to satisfy KKT, which by
Lemma~\ref{lem:kkt} requires
$a/(U/n) + \ell_i^\circ \equiv K^\star$ for all $i$, \ie $\ell_i^\circ$
constant.
\end{proof}

\paragraph{Edge case: small $U$.}
The KKT solution requires $K^\star > \max_i \ell_i^\circ$, equivalently
$U > n a/(K - \ell_{\max})$ for the relevant $K$. When $U$ is very
small, the optimum saturates: some $\sigma_i^2 \to 0$ for low-leverage
clients while the budget concentrates on high-leverage ones. This
regime falls below the deployment-relevant range and is not analysed
further.

\paragraph{Computational note.}
Solving \eqref{eq:budget-eq} numerically via 1-D bisection on $K$ is
$O(n \log(1/\epsilon))$ for tolerance $\epsilon$, trivial for $n \leq
500$. The implementation is in
\texttt{fulcrum/dp/allocation.py:optimal\_allocation}.
